\IfFileExists{stacks-project.cls}{%
\documentclass{stacks-project}
}{%
\documentclass{amsart}
}

% For dealing with references we use the comment environment
\usepackage{verbatim}
\newenvironment{reference}{\comment}{\endcomment}
%\newenvironment{reference}{}{}
\newenvironment{slogan}{\comment}{\endcomment}
\newenvironment{history}{\comment}{\endcomment}

% For commutative diagrams we use Xy-pic
\usepackage[all]{xy}

% We use 2cell for 2-commutative diagrams.
\xyoption{2cell}
\UseAllTwocells

% We use multicol for the list of chapters between chapters
\usepackage{multicol}

% This is generally recommended for better output
\usepackage{lmodern}
\usepackage[T1]{fontenc}

% For cross-file-references
\usepackage{xr-hyper}

% Package for hypertext links:
\usepackage{hyperref}

% For any local file, say "hello.tex" you want to link to please
% use \externaldocument[hello-]{hello}
\externaldocument[introduction-]{introduction}
\externaldocument[conventions-]{conventions}
\externaldocument[sets-]{sets}
\externaldocument[categories-]{categories}
\externaldocument[topology-]{topology}
\externaldocument[sheaves-]{sheaves}
\externaldocument[sites-]{sites}
\externaldocument[stacks-]{stacks}
\externaldocument[fields-]{fields}
\externaldocument[algebra-]{algebra}
\externaldocument[brauer-]{brauer}
\externaldocument[homology-]{homology}
\externaldocument[derived-]{derived}
\externaldocument[simplicial-]{simplicial}
\externaldocument[more-algebra-]{more-algebra}
\externaldocument[smoothing-]{smoothing}
\externaldocument[modules-]{modules}
\externaldocument[sites-modules-]{sites-modules}
\externaldocument[injectives-]{injectives}
\externaldocument[cohomology-]{cohomology}
\externaldocument[sites-cohomology-]{sites-cohomology}
\externaldocument[dga-]{dga}
\externaldocument[dpa-]{dpa}
\externaldocument[sdga-]{sdga}
\externaldocument[hypercovering-]{hypercovering}
\externaldocument[schemes-]{schemes}
\externaldocument[constructions-]{constructions}
\externaldocument[properties-]{properties}
\externaldocument[morphisms-]{morphisms}
\externaldocument[coherent-]{coherent}
\externaldocument[divisors-]{divisors}
\externaldocument[limits-]{limits}
\externaldocument[varieties-]{varieties}
\externaldocument[topologies-]{topologies}
\externaldocument[descent-]{descent}
\externaldocument[perfect-]{perfect}
\externaldocument[more-morphisms-]{more-morphisms}
\externaldocument[flat-]{flat}
\externaldocument[groupoids-]{groupoids}
\externaldocument[more-groupoids-]{more-groupoids}
\externaldocument[etale-]{etale}
\externaldocument[chow-]{chow}
\externaldocument[intersection-]{intersection}
\externaldocument[pic-]{pic}
\externaldocument[weil-]{weil}
\externaldocument[adequate-]{adequate}
\externaldocument[dualizing-]{dualizing}
\externaldocument[duality-]{duality}
\externaldocument[discriminant-]{discriminant}
\externaldocument[derham-]{derham}
\externaldocument[local-cohomology-]{local-cohomology}
\externaldocument[algebraization-]{algebraization}
\externaldocument[curves-]{curves}
\externaldocument[resolve-]{resolve}
\externaldocument[models-]{models}
\externaldocument[functors-]{functors}
\externaldocument[equiv-]{equiv}
\externaldocument[pione-]{pione}
\externaldocument[etale-cohomology-]{etale-cohomology}
\externaldocument[proetale-]{proetale}
\externaldocument[relative-cycles-]{relative-cycles}
\externaldocument[more-etale-]{more-etale}
\externaldocument[trace-]{trace}
\externaldocument[crystalline-]{crystalline}
\externaldocument[spaces-]{spaces}
\externaldocument[spaces-properties-]{spaces-properties}
\externaldocument[spaces-morphisms-]{spaces-morphisms}
\externaldocument[decent-spaces-]{decent-spaces}
\externaldocument[spaces-cohomology-]{spaces-cohomology}
\externaldocument[spaces-limits-]{spaces-limits}
\externaldocument[spaces-divisors-]{spaces-divisors}
\externaldocument[spaces-over-fields-]{spaces-over-fields}
\externaldocument[spaces-topologies-]{spaces-topologies}
\externaldocument[spaces-descent-]{spaces-descent}
\externaldocument[spaces-perfect-]{spaces-perfect}
\externaldocument[spaces-more-morphisms-]{spaces-more-morphisms}
\externaldocument[spaces-flat-]{spaces-flat}
\externaldocument[spaces-groupoids-]{spaces-groupoids}
\externaldocument[spaces-more-groupoids-]{spaces-more-groupoids}
\externaldocument[bootstrap-]{bootstrap}
\externaldocument[spaces-pushouts-]{spaces-pushouts}
\externaldocument[spaces-chow-]{spaces-chow}
\externaldocument[groupoids-quotients-]{groupoids-quotients}
\externaldocument[spaces-more-cohomology-]{spaces-more-cohomology}
\externaldocument[spaces-simplicial-]{spaces-simplicial}
\externaldocument[spaces-duality-]{spaces-duality}
\externaldocument[formal-spaces-]{formal-spaces}
\externaldocument[restricted-]{restricted}
\externaldocument[spaces-resolve-]{spaces-resolve}
\externaldocument[formal-defos-]{formal-defos}
\externaldocument[defos-]{defos}
\externaldocument[cotangent-]{cotangent}
\externaldocument[examples-defos-]{examples-defos}
\externaldocument[algebraic-]{algebraic}
\externaldocument[examples-stacks-]{examples-stacks}
\externaldocument[stacks-sheaves-]{stacks-sheaves}
\externaldocument[criteria-]{criteria}
\externaldocument[artin-]{artin}
\externaldocument[quot-]{quot}
\externaldocument[stacks-properties-]{stacks-properties}
\externaldocument[stacks-morphisms-]{stacks-morphisms}
\externaldocument[stacks-limits-]{stacks-limits}
\externaldocument[stacks-cohomology-]{stacks-cohomology}
\externaldocument[stacks-perfect-]{stacks-perfect}
\externaldocument[stacks-introduction-]{stacks-introduction}
\externaldocument[stacks-more-morphisms-]{stacks-more-morphisms}
\externaldocument[stacks-geometry-]{stacks-geometry}
\externaldocument[moduli-]{moduli}
\externaldocument[moduli-curves-]{moduli-curves}
\externaldocument[examples-]{examples}
\externaldocument[exercises-]{exercises}
\externaldocument[guide-]{guide}
\externaldocument[desirables-]{desirables}
\externaldocument[coding-]{coding}
\externaldocument[obsolete-]{obsolete}
\externaldocument[fdl-]{fdl}
\externaldocument[index-]{index}

% Theorem environments.
%
\theoremstyle{plain}
\newtheorem{theorem}[subsection]{Theorem}
\newtheorem{proposition}[subsection]{Proposition}
\newtheorem{lemma}[subsection]{Lemma}

\theoremstyle{definition}
\newtheorem{definition}[subsection]{Definition}
\newtheorem{example}[subsection]{Example}
\newtheorem{exercise}[subsection]{Exercise}
\newtheorem{situation}[subsection]{Situation}

\theoremstyle{remark}
\newtheorem{remark}[subsection]{Remark}
\newtheorem{remarks}[subsection]{Remarks}

\numberwithin{equation}{subsection}

% Macros
%
\def\lim{\mathop{\mathrm{lim}}\nolimits}
\def\colim{\mathop{\mathrm{colim}}\nolimits}
\def\Spec{\mathop{\mathrm{Spec}}}
\def\Hom{\mathop{\mathrm{Hom}}\nolimits}
\def\Ext{\mathop{\mathrm{Ext}}\nolimits}
\def\SheafHom{\mathop{\mathcal{H}\!\mathit{om}}\nolimits}
\def\SheafExt{\mathop{\mathcal{E}\!\mathit{xt}}\nolimits}
\def\Sch{\mathit{Sch}}
\def\Mor{\mathop{\mathrm{Mor}}\nolimits}
\def\Ob{\mathop{\mathrm{Ob}}\nolimits}
\def\Sh{\mathop{\mathit{Sh}}\nolimits}
\def\NL{\mathop{N\!L}\nolimits}
\def\CH{\mathop{\mathrm{CH}}\nolimits}
\def\proetale{{pro\text{-}\acute{e}tale}}
\def\etale{{\acute{e}tale}}
\def\QCoh{\mathit{QCoh}}
\def\Ker{\mathop{\mathrm{Ker}}}
\def\Im{\mathop{\mathrm{Im}}}
\def\Coker{\mathop{\mathrm{Coker}}}
\def\Coim{\mathop{\mathrm{Coim}}}

% Boxtimes
%
\DeclareMathSymbol{\boxtimes}{\mathbin}{AMSa}{"02}

%
% Macros for moduli stacks/spaces
%
\def\QCohstack{\mathcal{QC}\!\mathit{oh}}
\def\Cohstack{\mathcal{C}\!\mathit{oh}}
\def\Spacesstack{\mathcal{S}\!\mathit{paces}}
\def\Quotfunctor{\mathrm{Quot}}
\def\Hilbfunctor{\mathrm{Hilb}}
\def\Curvesstack{\mathcal{C}\!\mathit{urves}}
\def\Polarizedstack{\mathcal{P}\!\mathit{olarized}}
\def\Complexesstack{\mathcal{C}\!\mathit{omplexes}}
% \Pic is the operator that assigns to X its picard group, usage \Pic(X)
% \Picardstack_{X/B} denotes the Picard stack of X over B
% \Picardfunctor_{X/B} denotes the Picard functor of X over B
\def\Pic{\mathop{\mathrm{Pic}}\nolimits}
\def\Picardstack{\mathcal{P}\!\mathit{ic}}
\def\Picardfunctor{\mathrm{Pic}}
\def\Deformationcategory{\mathcal{D}\!\mathit{ef}}
\usepackage{cleveref}
\crefname{originallemma}{original lemma}{original lemmas}
\Crefname{originallemma}{Original lemma}{Original lemmas}
\makeatletter
\newlabel{varieties-lemma-ample-after-field-extension}{{0BDC}{source}{}{sourcepointer-1}{}}
\newlabel{varieties-lemma-ample-after-field-extension@cref}{{[originallemma][0][]0BDC}{[1][0][]source}{}{}{}}
\makeatother
\crefname{originalproposition}{original proposition}{original propositions}
\Crefname{originalproposition}{Original proposition}{Original propositions}
\makeatletter
\newlabel{proposition-connected-component}{{0B7R}{source}{}{sourcepointer-2}{}}
\newlabel{proposition-connected-component@cref}{{[originalproposition][0][]0B7R}{[1][0][]source}{}{}{}}
\makeatother
\crefname{originallemma}{original lemma}{original lemmas}
\Crefname{originallemma}{Original lemma}{Original lemmas}
\makeatletter
\newlabel{lemma-group-scheme-characteristic-zero-smooth}{{047N}{source}{}{sourcepointer-3}{}}
\newlabel{lemma-group-scheme-characteristic-zero-smooth@cref}{{[originallemma][0][]047N}{[1][0][]source}{}{}{}}
\makeatother
\crefname{originalproposition}{original proposition}{original propositions}
\Crefname{originalproposition}{Original proposition}{Original propositions}
\makeatletter
\newlabel{divisors-proposition-push-down-ample}{{0BD4}{source}{}{sourcepointer-4}{}}
\newlabel{divisors-proposition-push-down-ample@cref}{{[originalproposition][0][]0BD4}{[1][0][]source}{}{}{}}
\makeatother
\crefname{originallemma}{original lemma}{original lemmas}
\Crefname{originallemma}{Original lemma}{Original lemmas}
\makeatletter
\newlabel{more-morphisms-lemma-quasi-projective}{{0B42}{source}{}{sourcepointer-5}{}}
\newlabel{more-morphisms-lemma-quasi-projective@cref}{{[originallemma][0][]0B42}{[1][0][]source}{}{}{}}
\makeatother
\crefname{originallemma}{original lemma}{original lemmas}
\Crefname{originallemma}{Original lemma}{Original lemmas}
\makeatletter
\newlabel{lemma-reduced-subgroup-scheme-perfect}{{047R}{source}{}{sourcepointer-6}{}}
\newlabel{lemma-reduced-subgroup-scheme-perfect@cref}{{[originallemma][0][]047R}{[1][0][]source}{}{}{}}
\makeatother
\crefname{originallemma}{original lemma}{original lemmas}
\Crefname{originallemma}{Original lemma}{Original lemmas}
\makeatletter
\newlabel{lemma-reduced-group-scheme-prefect-field-characteristic-p-smooth}{{047P}{source}{}{sourcepointer-7}{}}
\newlabel{lemma-reduced-group-scheme-prefect-field-characteristic-p-smooth@cref}{{[originallemma][0][]047P}{[1][0][]source}{}{}{}}
\makeatother
\crefname{originallemma}{original lemma}{original lemmas}
\Crefname{originallemma}{Original lemma}{Original lemmas}
\makeatletter
\newlabel{varieties-lemma-smooth-regular}{{056S}{source}{}{sourcepointer-8}{}}
\newlabel{varieties-lemma-smooth-regular@cref}{{[originallemma][0][]056S}{[1][0][]source}{}{}{}}
\makeatother
\crefname{originallemma}{original lemma}{original lemmas}
\Crefname{originallemma}{Original lemma}{Original lemmas}
\makeatletter
\newlabel{more-algebra-lemma-regular-local-UFD}{{0AG0}{source}{}{sourcepointer-9}{}}
\newlabel{more-algebra-lemma-regular-local-UFD@cref}{{[originallemma][0][]0AG0}{[1][0][]source}{}{}{}}
\makeatother
\crefname{originallemma}{original lemma}{original lemmas}
\Crefname{originallemma}{Original lemma}{Original lemmas}
\makeatletter
\newlabel{divisors-lemma-complement-open-affine-effective-cartier-divisor}{{0BCW}{source}{}{sourcepointer-10}{}}
\newlabel{divisors-lemma-complement-open-affine-effective-cartier-divisor@cref}{{[originallemma][0][]0BCW}{[1][0][]source}{}{}{}}
\makeatother
\crefname{originallemma}{original lemma}{original lemmas}
\Crefname{originallemma}{Original lemma}{Original lemmas}
\makeatletter
\newlabel{lemma-group-scheme-over-field-separated}{{047L}{source}{}{sourcepointer-11}{}}
\newlabel{lemma-group-scheme-over-field-separated@cref}{{[originallemma][0][]047L}{[1][0][]source}{}{}{}}
\makeatother
\crefname{originallemma}{original lemma}{original lemmas}
\Crefname{originallemma}{Original lemma}{Original lemmas}
\makeatletter
\newlabel{morphisms-lemma-smooth-etale-over-affine-space}{{054L}{source}{}{sourcepointer-12}{}}
\newlabel{morphisms-lemma-smooth-etale-over-affine-space@cref}{{[originallemma][0][]054L}{[1][0][]source}{}{}{}}
\makeatother
\crefname{originalsection}{original section}{original sections}
\Crefname{originalsection}{Original section}{Original sections}
\makeatletter
\newlabel{divisors-section-norms}{{0BCX}{source}{}{sourcepointer-13}{}}
\newlabel{divisors-section-norms@cref}{{[originalsection][0][]0BCX}{[1][0][]source}{}{}{}}
\makeatother
\crefname{originallemma}{original lemma}{original lemmas}
\Crefname{originallemma}{Original lemma}{Original lemmas}
\makeatletter
\newlabel{divisors-lemma-weil-divisor-is-cartier-UFD}{{0AGA}{source}{}{sourcepointer-14}{}}
\newlabel{divisors-lemma-weil-divisor-is-cartier-UFD@cref}{{[originallemma][0][]0AGA}{[1][0][]source}{}{}{}}
\makeatother
\crefname{originaldefinition}{original definition}{original definitions}
\Crefname{originaldefinition}{Original definition}{Original definitions}
\makeatletter
\newlabel{divisors-definition-invertible-sheaf-effective-Cartier-divisor}{{01WX}{source}{}{sourcepointer-15}{}}
\newlabel{divisors-definition-invertible-sheaf-effective-Cartier-divisor@cref}{{[originaldefinition][0][]01WX}{[1][0][]source}{}{}{}}
\makeatother
\crefname{originallemma}{original lemma}{original lemmas}
\Crefname{originallemma}{Original lemma}{Original lemmas}
\makeatletter
\newlabel{varieties-lemma-rational-equivalence-for-Pic}{{0BEH}{source}{}{sourcepointer-16}{}}
\newlabel{varieties-lemma-rational-equivalence-for-Pic@cref}{{[originallemma][0][]0BEH}{[1][0][]source}{}{}{}}
\makeatother
\crefname{originaldefinition}{original definition}{original definitions}
\Crefname{originaldefinition}{Original definition}{Original definitions}
\makeatletter
\newlabel{properties-definition-ample}{{01PS}{source}{}{sourcepointer-17}{}}
\newlabel{properties-definition-ample@cref}{{[originaldefinition][0][]01PS}{[1][0][]source}{}{}{}}
\makeatother
\title{An algebraic group scheme over a field is quasi-projective}
\author{The Stacks Project Authors}
\date{stacks-project:groupoids.tex:local-snapshot-20261012; GNU FDL 1.2 or later}
\begin{document}
\maketitle
\setcounter{subsection}{1}
\noindent\textit{Editorial note.} All delivered mathematical text is the author's own, copied from the pinned Stacks Project chapter file groupoids.tex; no statement or proof step was rewritten, reordered or removed. The counted span is the author environment and its own proof environment. References inside the span are handed to the assembler as navigation stubs that carry the canonical tag of the cited tag only; whether each of them is load-bearing is recorded on the stub itself.\par\medskip
\noindent\textit{Allowed standard prerequisites (editorial).}\begin{itemize}
\item Lemma varieties-lemma-ample-after-field-extension (canonical tag 0BDC): ampleness descends from a field extension. cited at proof line 1588. Reduces to k algebraically closed.
\item Proposition proposition-connected-component (canonical tag 0B7R): the connected component G\textasciicircum{}0 of an algebraic group scheme over a field is an open and closed normal subgroup scheme whose cosets are isomorphic to G\textasciicircum{}0. cited at proof lines 1590-1592. Identifies the non-identity components and cuts them out of the argument.
\item Lemma lemma-group-scheme-characteristic-zero-smooth (canonical tag 047N) and Lemma lemma-reduced-group-scheme-prefect-field-characteristic-p-smooth (canonical tag 047P): smoothness in characteristic zero, resp. of the reduction in characteristic p; Lemma lemma-reduced-subgroup-scheme-perfect (canonical tag 047R) for H = G\_red being a group scheme. cited at proof lines 1598-1612. Reduces to the smooth irreducible case.
\item Divisors results: complement-of-affine-open is an effective Cartier divisor (0BCW), Weil divisors are Cartier over UFD local rings (0AGA), and the invertible sheaf of an effective Cartier divisor (01WX). cited at proof lines 1621-1671. Produce D on G and E on G x V and their invertible sheaves.
\item Lemma morphisms-lemma-smooth-etale-over-affine-space (canonical tag 054L). cited at proof lines 1634-1637. Supplies the etale cover U -> A\textasciicircum{}d\_k with the finite etale refinement W -> V.
\item Lemma varieties-lemma-rational-equivalence-for-Pic (canonical tag 0BEH). cited at proof lines 1663-1665. Shows the class L of the fibre divisors is independent of v, so the sections s\_v live in one invertible sheaf.
\item Definition properties-definition-ample (canonical tag 01PS). cited at proof lines 1676-1678. The exhibited sections satisfy the defining property of ampleness.
\end{itemize}\medskip
\ifdefined\hypertarget\hypertarget{sourcepointer-1}{}\fi
\noindent\textit{Source reference.} Original lemma 0BDC: sources/stacks-chow/tags.txt, source lines 14310--14310. Cited at proof line 1588 to reduce to the case of an algebraically closed ground field; load-bearing.\par\medskip
\ifdefined\hypertarget\hypertarget{sourcepointer-2}{}\fi
\noindent\textit{Source reference.} Original proposition 0B7R: sources/stacks-groupoids/pinned-original-source.tex, source lines 1230--1230. Cited at proof lines 1590-1592; the reduction to the identity component G\textasciicircum{}0 rests on it; load-bearing.\par\medskip
\ifdefined\hypertarget\hypertarget{sourcepointer-3}{}\fi
\noindent\textit{Source reference.} Original lemma 047N: sources/stacks-groupoids/pinned-original-source.tex, source lines 1483--1483. Cited at proof lines 1598-1600 for smoothness of G in characteristic zero; load-bearing.\par\medskip
\ifdefined\hypertarget\hypertarget{sourcepointer-4}{}\fi
\noindent\textit{Source reference.} Original proposition 0BD4: sources/stacks-chow/tags.txt, source lines 14302--14302. Cited at proof lines 1602-1604: it suffices to exhibit an ample invertible sheaf on the reduction H; load-bearing.\par\medskip
\ifdefined\hypertarget\hypertarget{sourcepointer-5}{}\fi
\noindent\textit{Source reference.} Original lemma 0B42: sources/stacks-chow/tags.txt, source lines 13985--13985. Cited parenthetically at proof lines 1604-1607 to identify having an ample invertible sheaf with quasi-projectivity; navigation for the framing.\par\medskip
\ifdefined\hypertarget\hypertarget{sourcepointer-6}{}\fi
\noindent\textit{Source reference.} Original lemma 047R: sources/stacks-groupoids/pinned-original-source.tex, source lines 1063--1063. Cited at proof lines 1609-1610 to know H = G\_red is a group scheme; load-bearing.\par\medskip
\ifdefined\hypertarget\hypertarget{sourcepointer-7}{}\fi
\noindent\textit{Source reference.} Original lemma 047P: sources/stacks-groupoids/pinned-original-source.tex, source lines 1512--1512. Cited at proof lines 1610-1612 for smoothness of H in characteristic p; load-bearing.\par\medskip
\ifdefined\hypertarget\hypertarget{sourcepointer-8}{}\fi
\noindent\textit{Source reference.} Original lemma 056S: sources/stacks-chow/tags.txt, source lines 6720--6720. Cited at proof lines 1617-1619 to see the local rings of G are regular; load-bearing.\par\medskip
\ifdefined\hypertarget\hypertarget{sourcepointer-9}{}\fi
\noindent\textit{Source reference.} Original lemma 0AG0: sources/stacks-chow/tags.txt, source lines 13178--13178. Cited at proof lines 1618-1620 to upgrade regular local rings to UFDs; load-bearing.\par\medskip
\ifdefined\hypertarget\hypertarget{sourcepointer-10}{}\fi
\noindent\textit{Source reference.} Original lemma 0BCW: sources/stacks-chow/tags.txt, source lines 14294--14294. Cited at proof lines 1621-1626 to produce the Cartier divisor D with G \textbackslash{}\textbackslash{} D affine; load-bearing.\par\medskip
\ifdefined\hypertarget\hypertarget{sourcepointer-11}{}\fi
\noindent\textit{Source reference.} Original lemma 047L: sources/stacks-groupoids/pinned-original-source.tex, source lines 983--983. Cited parenthetically at proof lines 1624-1625; used when intersecting affine opens inside G; load-bearing.\par\medskip
\ifdefined\hypertarget\hypertarget{sourcepointer-12}{}\fi
\noindent\textit{Source reference.} Original lemma 054L: sources/stacks-chow/tags.txt, source lines 6644--6644. Cited at proof lines 1634-1637 to choose the etale cover U -> A\textasciicircum{}d\_k; load-bearing.\par\medskip
\ifdefined\hypertarget\hypertarget{sourcepointer-13}{}\fi
\noindent\textit{Source reference.} Original section 0BCX: sources/stacks-chow/tags.txt, source lines 14295--14295. Cited inside the footnote at proof lines 1638-1641 as an alternative route to the divisor E; navigation for an alternative, not for the counted argument.\par\medskip
\ifdefined\hypertarget\hypertarget{sourcepointer-14}{}\fi
\noindent\textit{Source reference.} Original lemma 0AGA: sources/stacks-chow/tags.txt, source lines 13188--13188. Cited at proof lines 1645-1649 to conclude the components of T are effective Cartier divisors; load-bearing.\par\medskip
\ifdefined\hypertarget\hypertarget{sourcepointer-15}{}\fi
\noindent\textit{Source reference.} Original definition 01WX: sources/stacks-chow/tags.txt, source lines 2498--2498. Cited at proof line 1662 to define M = O(E) and at lines 1668-1671 for the canonical section; load-bearing.\par\medskip
\ifdefined\hypertarget\hypertarget{sourcepointer-16}{}\fi
\noindent\textit{Source reference.} Original lemma 0BEH: sources/stacks-chow/tags.txt, source lines 14350--14350. Cited at proof lines 1663-1665 to see the isomorphism class L of M\_v is independent of v; load-bearing.\par\medskip
\ifdefined\hypertarget\hypertarget{sourcepointer-17}{}\fi
\noindent\textit{Source reference.} Original definition 01PS: sources/stacks-chow/tags.txt, source lines 2241--2241. Cited at proof lines 1676-1678 to conclude L is ample by definition; load-bearing.\par\medskip
\begin{lemma}
\label{lemma-algebraic-quasi-projective}
Let $k$ be a field. Let $G$ be an algebraic group scheme over $k$.
Then $G$ is quasi-projective over $k$.
\end{lemma}

\begin{proof}
By Varieties, Lemma \ref{varieties-lemma-ample-after-field-extension}
we may assume that $k$ is algebraically closed. Let $G^0 \subset G$
be the connected component of $G$ as in
Proposition \ref{proposition-connected-component}.
Then every other connected component of $G$ has a $k$-rational
point and hence is isomorphic to $G^0$ as a scheme.
Since $G$ is quasi-compact and Noetherian, there are finitely many of these
connected components. Thus we reduce to the case discussed in
the next paragraph.

\medskip\noindent
Let $G$ be a connected algebraic group scheme over an algebraically closed
field $k$. If the characteristic of $k$ is zero, then $G$ is smooth over
$k$ by Lemma \ref{lemma-group-scheme-characteristic-zero-smooth}.
If the characteristic of $k$ is $p > 0$, then we let $H = G_{red}$
be the reduction of $G$. By
Divisors, Proposition \ref{divisors-proposition-push-down-ample}
it suffices to show that $H$ has an ample invertible sheaf.
(For an algebraic scheme over $k$ having an ample invertible
sheaf is equivalent to being quasi-projective over $k$, see
for example the very general
More on Morphisms, Lemma \ref{more-morphisms-lemma-quasi-projective}.)
By Lemma \ref{lemma-reduced-subgroup-scheme-perfect}
we see that $H$ is a group scheme over $k$.
By Lemma \ref{lemma-reduced-group-scheme-prefect-field-characteristic-p-smooth}
we see that $H$ is smooth over $k$.
This reduces us to the situation discussed in the next
paragraph.

\medskip\noindent
Let $G$ be a quasi-compact irreducible smooth group scheme over an
algebraically closed field $k$. Observe that the local rings of $G$
are regular and hence UFDs
(Varieties, Lemma \ref{varieties-lemma-smooth-regular} and
More on Algebra, Lemma \ref{more-algebra-lemma-regular-local-UFD}).
The complement of a nonempty affine open of $G$
is the support of an effective Cartier divisor $D$.
This follows from Divisors, Lemma
\ref{divisors-lemma-complement-open-affine-effective-cartier-divisor}.
(Observe that $G$ is separated by
Lemma \ref{lemma-group-scheme-over-field-separated}.)
We conclude there exists an effective Cartier divisor $D \subset G$
such that $G \setminus D$ is affine. We will use below that
for any $n \geq 1$ and $g_1, \ldots, g_n \in G(k)$ the complement
$G \setminus \bigcup D g_i$ is affine. Namely, it is the intersection
of the affine opens $G \setminus Dg_i \cong G \setminus D$
in the separated scheme $G$.

\medskip\noindent
We may choose the top row of the diagram
$$
\xymatrix{
G & U \ar[l]_j \ar[r]^\pi & \mathbf{A}^d_k \\
& W \ar[r]^{\pi'} \ar[u] & V \ar[u]
}
$$
such that $U \not = \emptyset$, $j : U \to G$ is an open immersion, and
$\pi$ is \'etale, see
Morphisms, Lemma \ref{morphisms-lemma-smooth-etale-over-affine-space}.
There is a nonempty affine open $V \subset \mathbf{A}^d_k$ such that
with $W = \pi^{-1}(V)$ the morphism $\pi' = \pi|_W : W \to V$ is finite \'etale.
In particular $\pi'$ is finite locally free, say of degree $n$.
Consider the effective Cartier divisor
$$
\mathcal{D} = \{(g, w) \mid m(g, j(w)) \in D\} \subset G \times W
$$
(This is the restriction to $G \times W$ of the pullback of $D \subset G$
under the flat morphism $m : G \times G \to G$.)
Consider the closed subset\footnote{Using the material
in Divisors, Section \ref{divisors-section-norms}
we could take as effective Cartier
divisor $E$ the norm of the effective Cartier divisor $\mathcal{D}$
along the finite locally free morphism $1 \times \pi'$ bypassing
some of the arguments.}
$T = (1 \times \pi')(\mathcal{D}) \subset G \times V$.
Since $\pi'$ is finite locally free, every irreducible component
of $T$ has codimension $1$ in $G \times V$. Since $G \times V$
is smooth over $k$ we conclude these components are effective Cartier
divisors (Divisors, Lemma \ref{divisors-lemma-weil-divisor-is-cartier-UFD}
and lemmas cited above)
and hence $T$ is the support of an effective Cartier divisor
$E$ in $G \times V$. If $v \in V(k)$, then
$(\pi')^{-1}(v) = \{w_1, \ldots, w_n\} \subset W(k)$ and we see that
$$
E_v = \bigcup\nolimits_{i = 1, \ldots, n} D j(w_i)^{-1}
$$
in $G$ set theoretically. In particular we see that $G \setminus E_v$
is affine open (see above).
Moreover, if $g \in G(k)$, then there exists a $v \in V$
such that $g \not \in E_v$. Namely, the set $W'$ of $w \in W$ such that
$g \not \in Dj(w)^{-1}$ is nonempty open and it suffices to pick $v$
such that the fibre of $W' \to V$ over $v$ has $n$ elements.

\medskip\noindent
Consider the invertible sheaf
$\mathcal{M} = \mathcal{O}_{G \times V}(E)$ on $G \times V$.
By Varieties, Lemma \ref{varieties-lemma-rational-equivalence-for-Pic}
the isomorphism class $\mathcal{L}$ of the restriction
$\mathcal{M}_v = \mathcal{O}_G(E_v)$ is independent of $v \in V(k)$.
On the other hand, for every $g \in G(k)$ we can find a $v$
such that $g \not \in E_v$ and such that $G \setminus E_v$
is affine. Thus the canonical section
(Divisors, Definition
\ref{divisors-definition-invertible-sheaf-effective-Cartier-divisor})
of $\mathcal{O}_G(E_v)$
corresponds to a section $s_v$ of $\mathcal{L}$ which does not
vanish at $g$ and such that $G_{s_v}$ is affine.
This means that $\mathcal{L}$ is ample by definition
(Properties, Definition \ref{properties-definition-ample}).
\end{proof}

\end{document}
